\documentclass[10pt,a4paper]{amsart}
\usepackage[margin=19mm]{geometry}
\usepackage{amsmath,amssymb,mathtools}
\usepackage[scr=rsfs,cal=euler]{mathalfa}
\usepackage{xfrac}
\usepackage[unicode,hidelinks,bookmarks=false]{hyperref}
\newcommand{\ie}{i.e.\ }
\newcommand{\cf}{cf.\ }
\newcommand{\resp}{respectively\ }
\newcommand{\Subsection}{\subsection}
\providecommand{\nobreakdash}{\nobreak\hbox{-}}
\setlength{\emergencystretch}{2em}
\multlinegap0pt
\usepackage{amssymb}
\usepackage{xcolor}
\newtheorem{example}{Example}
\newtheorem{lemma}{Lemma}
\def\Endoftheorem{\@endgadget{\@Endofsymbol}}
\newcommand{\T}{\textsc{t}}
\newcommand{\e}{\hspace{.4pt}\mathrm{e}}
\newcommand{\ep}{\odot}
\newcommand{\mC}{\mathbb{C}}
\newcommand{\mR}{\mathbb{R}}
\newcommand{\mi}{\hspace{.4pt}\mathrm{i}}
\newcommand{\mymatrix}[1]{\begin{bmatrix} #1 \end{bmatrix}}
\newcommand{\smat}[1]{\bigl[\begin{smallmatrix} #1 
                            \end{smallmatrix}\bigr]} 
\newcommand{\ssmat}[1]{\begin{smallmatrix} #1 \end{smallmatrix}}
\title{Planar Pendulums in Complex Coordinates}
\author{Gjerrit Meinsma}
\date{Source: arXiv:2609.15448v1 (CC BY 4.0)}
\begin{document}
\maketitle

\noindent\emph{Source and licence.} Planar Pendulums in Complex Coordinates. Version arXiv:2609.15448v1 (CC BY 4.0), distributed by arXiv under the Creative Commons Attribution 4.0 International licence, http://creativecommons.org/licenses/by/4.0/, as recorded in the arXiv licence metadata for this exact version and shown on the abstract page https://arxiv.org/abs/2609.15448v1. Full licence notice: \texttt{licenses/arxiv-2609.15448-CC-BY-4.0.txt}.\par\smallskip

\noindent\textit{Editorial scope.} Counted unit: the article's lemma lem:dual (source0.tex lines 403--408). The article's complete original proof of it is delivered with it (source0.tex lines 409--462, one \texttt{proof} environment of 54 lines). No mathematics is written, repaired or re-derived: the statement and the proof are the article's own lines, in the article's own notation, and no retained line is substituted or re-flowed.  Retained uncounted as the source-local context the proof needs: lines 292--299; lines 304--348; lines 377--380; lines 395--397; lines 463--470.  One declared adaptation: exactly 2 ancillary strings inside the retained spans are omitted (image-inclusion commands and, where the source repeats a label, the repeated label definitions) and no other line differs from the source; the figure environments, with their placement, captions and labels, are kept, and the package records the omitted strings in fidelity receipts bound to the affected bodies.  No retained line contains a citation, so the wrapper carries no bibliography and no bibliography entry.  The retained lines contain the article's own diagram code, which the wrapper typesets with the standard tikz packages; no image file is delivered and the package declares no asset dependency.  Editorial formatting only, in addition: an AMS \texttt{amsart} preamble that restates the article's own declarations and macros used by the retained lines, namely the environments \texttt{example}, \texttt{lemma} on separate counters, together with the standard packages those lines need.  The retained environments print in this document as Example 1, Lemma 1 in order of appearance, whereas the article prints the same items with the numbering of its own full document.  The complete native source of the article as distributed in the arXiv e-print for this exact version is bundled unchanged as \texttt{sources/210416/source0.tex}.
\begin{figure}[htbp]
  \centering
  
  \caption{Single-arm pendulum in $\mC$ (see 
    Example~\ref{SP}). The compression forces are shown in red and
    blue. The gravitational force is in shown in orange.}
  \label{fig:SDP}
\end{figure}

\begin{example}[Single-arm pendulum]\label{SP} Consider the
  single-arm pendulum in $\mC$ as depicted in Fig.~\ref{fig:SDP}. The
  pendulum (its one arm) can rotate freely around the pivot at
  $0\in\mC$.  Let $\ell$ be the length of the arm, and let $z\in\mC$
  be the position of the tip of the pendulum. Hence,
  \[
    z=\ell\e^{\mi\phi},
  \]
  where ``$\mi$'' is the imaginary unit, and $\phi$ is the argument of the
  complex number $z$. The speed and acceleration,
  $\dot{z},\ddot{z}\in\mC$, thus are
  \begin{align*}
    \dot{z} 
    = \ell\e^{\mi\phi}\mi\dot{\phi},\qquad
    \ddot{z}
    = \ell\e^{\mi\phi}[\mi\ddot{\phi}-\dot{\phi}^2].
  \end{align*}
  The force $F\in\mC$ acting on the tip of the arm (at $z$) consists
  of gravity $gm$ (downwards, so in the ``$-\mi$ direction''), and a
  compression force $\lambda$ in the direction of the pendulum, see
  Fig.~\ref{fig:SDP}. Hence the total force acting on the tip is
  $-\mi g m +\lambda \e^{\mi\phi}\in\mC$ for some as yet unknown
  compression force $\lambda\in\mR$. Newton's 2nd law, $m\ddot{z}=F$,
  therefore becomes
  \[
    m\ell\e^{\mi\phi}[\mi\ddot{\phi}-\dot{\phi}^2]
    = -\mi g m +\lambda\e^{\mi\phi}. 
  \]
  Dividing by $m\ell\e^{\mi\phi}$ gives
  \[
    \mi\ddot{\phi}-\dot{\phi}^2 = -\mi
    \frac{g}{\ell}\e^{-\mi\phi}+\frac{1}{m\ell}\lambda.
  \]
  This equation is linear in the real pair $(\ddot{\phi},\lambda)$,
  and its imaginary and real part respectively determine $\ddot{\phi}$
  and $\lambda$:
  \begin{align*}
    \ddot{\phi} &=-\frac{g}{\ell}\cos(\phi),\\
    \lambda & = -m\ell \dot{\phi}^2+gm \sin(\phi). 
  \end{align*}
  The first of these is the equation of motion. Once we have
  $\phi,\dot{\phi}$, the second equation determines the compression
  force $\lambda$.
  \Endoftheorem
\end{example}

\begin{equation}
  \label{eq:cm}
  z= S[\ell\ep\e^{\mi\phi}]
\end{equation}

\begin{equation}\label{fizzz}
  F=-\mi g m+ B[\lambda\ep\e^{\mi\phi}]
\end{equation}

\begin{figure}
  \centering
  
  \caption{An arbitrary pendulum with $n+1$ arms, seen as a pendulum
    with $n$ arms (light gray) with one more arm added (dark gray).
    The compression force in the added arm is indicated in red and blue.}
  \label{fig:fant}
\end{figure}

\begin{lemma}[Duality]\label{lem:dual}
  For every pendulum tree with $n$ arms we have that~\eqref{eq:cm}
  and~\eqref{fizzz} hold for some $S\in\{0,1\}^{n\times n}$ and
  $B\in\{-1,0,1\}^{n\times n}$. In fact, $B=S^{-\T}$, and all
  eigenvalues of $S$ and $B$ are equal to 1.
\end{lemma}

\begin{proof}
  We prove it by induction in the number of arms $n$. For $n=1$ arms
  (the single-arm pendulum) we have $S=B=1$ (see Example~\ref{SP}) so
  then the result holds.

  Next let $n\in\mathbb{N}$ and suppose that $B=S^{-\T}$ for every
  pendulum tree of $n$ arms, and that all eigenvalues of $S$ and $B$
  are $1$, and $S\in\{0,1\}^{n\times n}$ and
  $B\in\{-1,0,1\}^{n\times n}$ (the induction hypothesis).  Now
  consider an arbitrary pendulum tree with $n+1$ arms. By removing one
  arm it reduces to an $n$-arm pendulum. Without loss of generality we
  denote the tip of the removed arm by $z_{n+1}$, and the tip to which
  it was connected by $z_n$, see Fig.~\ref{fig:fant}. Thus,
  $z_{n+1}=z_n+\ell_{n+1}\e^{\mi\phi_{n+1}}$. The connectivity matrix
  $S_{n+1}$ of the $n+1$-arm pendulum in terms of the connectivity
  matrix $S$ of the $n$-arm pendulum follows from
  \[
    \mymatrix{z\\ z_{n+1}}
    =
    \mymatrix{z\\ z_n + \ell_{n+1}\e^{\mi\phi_{n+1}}}
    = \underbrace{\mymatrix{S & 0\\ S_{n*} &
        1}}_{S_{n+1}}\mymatrix{\ell\e^{\mi\phi}\\ \ell_{n+1}\e^{\mi\phi_{n+1}}},
  \]
  where $S_{n*}$ means the final row of $S$. By adding arm number
  $n+1$ to the pendulum, the total force $\smat{F\\ F_{n+1}}$ (acting
  on the $n+1$ joints/tips) becomes
  \[
    \mymatrix{F\\ F_{n+1}}
    = -\mi g \mymatrix{m\\ m_{n+1}}
    +\underbrace{\left[\begin{array}{c|c}
                         B & \ssmat{0\\[-1ex] \vdots\\ 0\\-1}\\ \hline
                         0 & 1
                       \end{array}\right]}_{B_{n+1}}
                   \mymatrix{\lambda\e^{\mi\phi}\\
                     \lambda_{n+1}\e^{\mi\phi_{n+1}}},
  \]
  (this defines $B_{n+1}$).  This equation expresses that there is one
  additional compression force $\lambda_{n+1}\e^{\mi\phi_{n+1}}$ and
  that it acts on $z_{n+1}$, and also on $z_n$ but then in opposite direction
  (see the red and blue forces in Fig.~\ref{fig:fant}).  By the
  induction hypothesis we have $B=S^{-\T}$, so
  \[
    S_{n+1}^{\T}B_{n+1}=
    \left[\begin{array}{c|c} S^{\T} & S_{n*}^{\T} \\ \hline 0 & 1\end{array}\right]
    \left[\begin{array}{c|c} S^{-\T} &
      \ssmat{0\\[-1ex] \vdots\\ 0\\-1}\\ \hline 0 & 1\end{array}\right]
      =\left[\begin{array}{c|c} I & 0 \\
           \hline 0 & 1\end{array}\right]=I_{n+1}.
  \]
  Hence, $B_{n+1}=S_{n+1}^{-\T}$, and the eigenvalues of $S_{n+1}$ are
  those of $S$ with one more eigenvalue equal to 1. So, by the
  induction hypothesis, they are all equal to $1$. Also, we see that
  $S\in\{0,1\}^{n\times n}$ and $B\in\{-1,0,1\}^{n\times n}$.
\end{proof}

\end{document}
